% ============================================================
%  Ishaq Hamza: Academic Curriculum Vitae
%  Compile with pdfLaTeX (e.g. on Overleaf).
% ============================================================
\documentclass[11pt,letterpaper]{article}

\usepackage[T1]{fontenc}
\usepackage{lmodern}
\usepackage[utf8]{inputenc}
\usepackage[letterpaper,margin=1in,top=0.75in,bottom=0.8in]{geometry}
\usepackage{titlesec}
\usepackage[shortlabels]{enumitem}
\usepackage{xcolor}
\usepackage{microtype}
\usepackage{amsmath}
\usepackage{multicol}
\usepackage{fancyhdr}
\usepackage{lastpage}
\usepackage{needspace}
\usepackage{hyperref}

\definecolor{linkblue}{RGB}{24,62,150}
\definecolor{rulegray}{gray}{0.55}
\definecolor{softgray}{gray}{0.35}
\hypersetup{
  colorlinks=true, urlcolor=linkblue, linkcolor=linkblue,
  pdftitle={Ishaq Hamza - Curriculum Vitae}, pdfauthor={Ishaq Hamza}
}

\setlength{\parindent}{0pt}
\setlength{\parskip}{0pt}
\raggedbottom

% ---------- Page footer ----------
\pagestyle{fancy}
\fancyhf{}
\renewcommand{\headrulewidth}{0pt}
\fancyfoot[L]{\small\color{softgray}Ishaq Hamza \textendash{} Curriculum Vitae}
\fancyfoot[R]{\small\color{softgray}Page \thepage\ of \pageref*{LastPage}}
\fancypagestyle{plain}{\fancyhf{}\fancyfoot[R]{\small\color{softgray}Page \thepage\ of \pageref*{LastPage}}}

% ---------- Section headings: small caps with a hairline rule ----------
\titleformat{\section}{\Large\scshape\raggedright}{}{0em}{}[{\color{rulegray}\vspace{1pt}\titlerule}]
\titlespacing*{\section}{0pt}{15pt}{8pt}

% ---------- Lists ----------
\setlist[itemize]{before={\interlinepenalty=10000}, leftmargin=1.2em, labelsep=0.5em, itemsep=2pt, topsep=3pt, parsep=0pt, partopsep=0pt}
\setlist[enumerate]{leftmargin=2.2em, itemsep=2pt, topsep=3pt, parsep=0pt}

% ---------- Entry helpers ----------
% \entry{Title}{Dates}
\newcommand{\entry}[2]{\needspace{6\baselineskip}\textbf{#1}\hfill #2\par}
% \sub{left (italic)}{right}
\newcommand{\sub}[2]{{\itshape #1}\hfill{#2}\par}
% \award{Title}{Year}{Description}
\newcommand{\award}[3]{\item \textbf{#1}\hfill #2\\ #3}
\newcommand{\gap}{\vspace{8pt}}

\begin{document}

% ======================= HEADER =======================
\begin{center}
  {\fontsize{25}{30}\selectfont\scshape Ishaq Hamza}\par\vspace{5pt}
  PhD Student in Operations Research\par
  Edwardson School of Industrial Engineering, Purdue University\par\vspace{4pt}
  \href{mailto:ihamza@purdue.edu}{ihamza@purdue.edu}
  \enspace$\vert$\enspace
  \href{https://scholar.google.com/citations?user=RCmlLGgAAAAJ&hl=en}{Google Scholar}
  \enspace$\vert$\enspace
  \href{https://www.linkedin.com/in/ishaqhamza}{LinkedIn}
  \enspace$\vert$\enspace
  \href{https://github.com/IshaqHamza}{GitHub}
\end{center}

% ======================= RESEARCH INTERESTS =======================
\section*{Research Interests}
My research centers on reinforcement learning theory and data-driven decision-making under uncertainty.
I design and analyze practical policy-space algorithms, and establish
finite-sample guarantees for them under minimal and realistic assumptions, aiming for theory that gives
practitioners real insight. More broadly, I am interested in stochastic
approximation, stochastic games, and in developing reinforcement learning into
a tool that reaches beyond its classical applications, assumptions, and modeling choices.

% ======================= EDUCATION =======================
\section*{Education}
\entry{Purdue University}{2026 \textendash{} Present}
\sub{Ph.D. in Operations Research, Edwardson School of Industrial Engineering}{}
Advisor: \href{https://sites.google.com/view/zaiweichen/home}{Dr.~Zaiwei Chen}\par

\gap
\entry{Indian Institute of Science (IISc)}{2022 \textendash{} 2026}
\sub{B.Tech in Math and Computing}{}
CGPA: \textbf{3.96/4}, \textbf{Rank 4}\par
\vspace{5pt}
Relevant coursework:
\begin{multicols}{2}
\begin{itemize}[itemsep=0pt, topsep=1pt]
  \item On-line and Reinforcement Learning
  \item Game Theory
  \item Measure Theoretic Probability
  \item Regression and Time-Series Analysis
  \item Stochastic Finance
  \item Applied DS\&AI
  \item Renewal and Queuing Theory
  \item Computational Optimization
  \item Detection \& Estimation Theory
  \item Numerical Methods
  \item Nonlinear Systems \& Stability
  \item Advanced Real Analysis
\end{itemize}
\end{multicols}

% ======================= AWARDS =======================
\section*{Awards and Achievements}
\begin{itemize}[itemsep=4pt]
  \award{Dr.~W. Dale and Mrs.~Jeanne C. Compton Award, Purdue University}{2026 \textendash{} 2030}
        {Four-year doctoral fellowship, awarded over and above standard PhD support.}
  \award{PhD Admission Offers}{2026}
        {Georgia Tech (ISyE) and Columbia University (IEOR).}
  \award{ITTIAM Fellowship, IISc}{2024\textendash25, 2025\textendash26}
        {Two-time awardee.
        \href{https://odaa.iisc.ac.in/ittiam-btech-fellowship/}{[details]}}
  \award{Dr.~Kolar \& Mrs.~Girija Kodandapani Scholarship, IISc}{}
        {Awarded to one student per batch.
        \href{https://odaa.iisc.ac.in/kodandapani-family-fellowship/}{[details]}}
  \award{Simon Marais Mathematics Competition}{2024}
        {Global rank 98 (5th within IISc).
        \href{https://www.simonmarais.org/east-division1.html}{[results]}}
  \award{Indian Statistical Institute Entrance Test}{2022}
        {National rank 62.}
\end{itemize}

% ======================= PUBLICATIONS =======================
\section*{Publications}
\begin{enumerate}[label={[\arabic*]}]
  \item Achieving $\epsilon^{-2}$ Sample Complexity in Single-Loop Actor-Critic,
        \href{https://openreview.net/forum?id=ALtp3rFOx7}{NeurIPS MLxOR 2026}, \textbf{Ishaq~Hamza}, Zaiwei~Chen.
        Undergrad publication, \href{https://arxiv.org/abs/2605.13639}{arXiv link}.
\end{enumerate}

% ======================= PRESENTATIONS =======================
\section*{Presentations}
\textit{To be updated.}

% ======================= POSITIONS =======================
\section*{Positions}
\entry{Summer Research Intern}{May \textendash{} Aug 2025}
\sub{Edwardson School of Industrial Engineering, Purdue University}{Advisor: Dr.~Zaiwei Chen}
Project: \textit{Optimal Single-Timescale Actor-Critic Algorithm}
\begin{itemize}
  \item Proposed two variants of the actor-critic algorithm with single-timescale updates, and showed through
        a tight analysis that they achieve the optimal sample complexity in the off-policy setting under
        provably minimal assumptions.
  \item The analysis combines Moreau envelopes, stochastic approximation, and Markov chain mixing.
        This work became publication [1].
\end{itemize}

\gap
\entry{Teaching Assistant, Advanced Convex Optimization}{Jan \textendash{} Apr 2025}
\sub{Department of Electrical Engineering, IISc}{Instructor: Dr.~Kunal N. Chaudhury}
A graduate course for PhD students, assisted during my third undergraduate year.
\href{https://sites.google.com/site/kunalnchaudhury/Teaching/aco-2025}{[course webpage]}
\begin{itemize}
  \item Conducted tutorials, answered students' questions, and graded assignments.
  \item The course covered rigorous convex analysis in extended-real settings and the operator-theoretic
        treatment of optimization algorithms in both smooth and non-smooth settings, using set-valued
        analysis for proximal algorithms.
\end{itemize}

% ======================= PROJECTS =======================
\section*{Pre-PhD Projects}
\entry{Extending Soft Actor-Critic to Multi-Agent RL}{Spring 2025}
\sub{RL course project}{\href{https://github.com/IshaqHamza/Multi-agent-SAC}{[GitHub]}}
\begin{itemize}
  \item Surveyed entropy-regularized RL and Soft Actor-Critic algorithms in several settings.
  \item Proposed and implemented a novel multi-agent extension of the algorithm, incorporating
        cross-entropy regularization.
  \item Extended the analysis to prove single-step policy evaluation and improvement theorems, and gave a
        counterexample to global optimal convergence.
  \item Proposed dynamic cross-entropy regularization to escape sub-optimal Nash equilibria and designed
        experiments to validate it; the rest of the team implemented them, with positive results.
\end{itemize}

\gap
\entry{Game-Theoretic Market-Making and Trading}{Spring 2025}
\sub{Game Theory course project}{}
\begin{itemize}
  \item Surveyed 10+ papers on the modeling and equilibrium analysis of market makers and traders in
        financial markets, and 5+ papers on information asymmetry and its impact on market making.
  \item Formulated a novel generalized model of trading with priced information exchange.
\end{itemize}

\gap
\entry{Automated Equity Research Analyst}{Fall 2024}
\sub{Applied DS\&AI course project (team of 6)}{\href{https://github.com/Finance-Automation-Projects/finance-dashboard}{[GitHub]}}
\begin{itemize}
  \item Conceived the rating-generation mechanism and the project structure for the team.
  \item Implemented a self-engineered variant of the multiplicative weights algorithm that rates stocks from
        aspect-based sentiment analysis scores computed on financial news.
  \item Tested it on both real and geometric-Brownian-motion-simulated prices, with high-noise expert ratings.
\end{itemize}

\gap
\entry{Financial Network Modeling}{Fall 2024}
\sub{Data Analytics course project (led a team of 5)}{\href{https://github.com/IshaqHamza/DA_Financial_networks}{[GitHub]}}
\begin{itemize}
  \item Studied generative models for analyzing community structure in graphs.
  \item Modeled relationships between stock price movements with a hierarchical weighted stochastic block model.
  \item Applied the same model to import--export data across 200+ countries to study trade relations,
        obtaining the expected community structure.
\end{itemize}

\vfill
\begin{flushright}\small\color{softgray}Last updated: October 2026\end{flushright}

\end{document}
